\documentclass[11pt]{article}
\usepackage[utf8]{inputenc}
\usepackage{amsmath,amssymb}
\usepackage{graphicx}
\usepackage{booktabs}
\usepackage{authblk}
\usepackage[margin=1in]{geometry}
\usepackage[hidelinks]{hyperref}
\usepackage{setspace}

\newcommand{\grad}{\nabla}
\newcommand{\half}{\tfrac12}


\title{RealQM: Insulation vs Conduction\\ \large the Free Boundary as Carrier}

\author[1]{Claes Johnson}
\affil[1]{KTH Royal Institute of Technology, SE-100 44 Stockholm, Sweden \\
\texttt{claesjohnson@gmail.com}}
\date{\today \\[0.4em]
\small\textit{Mathematical formulation by the author; analysis developed} \\
\small\textit{with AI assistance (Claude, Anthropic).}}

\begin{document}
\maketitle

\begin{abstract}
\noindent \textbf{Can RealQM explain the difference between a conductor and an insulator?} The
established answer is band theory: fill the bands of a periodic potential and count --- correct,
cheap, and applicable at the size of a crystal, but it supplies the potential rather than deriving
it, as density functional theory supplies a functional. RealQM starts further back, from $N$
non-negative unit charge densities in ordinary three-dimensional space, one per electron, meeting at
free boundaries and minimising the Coulomb energy, with nothing fitted and polynomial cost. This
article asks whether that reaches the same distinction, and where it stops.

\textbf{Polarisation, measured; sliding, not found.} On a regular chain of atoms with one valence
electron each, the polarisability $\alpha$ is extensive in the number of electrons: $\alpha=27.4N$,
flat to $1\%$ per atom --- the response of charge fixed per site, not of charge free to move between
sites. And no sliding of the partition could be identified: the corrugation it would have to pass
through vanishes at no spacing, filling or boundary treatment tested, placing the threshold field ---
for the mode measured, a rigid slide on a clamped lattice --- beyond dielectric breakdown. Both point
to insulation.

\textbf{The carrier, identified.} What moves is not a particle but a \emph{free boundary}, in the
relation a dislocation bears to slip --- which places protonic and ionic conduction inside the
framework and the electronic \emph{rate} --- a conductivity $\sigma$ --- outside it, in the
time-dependent form: a minimisation has no time in it.
\end{abstract}

\setstretch{1.15}

\section{The question, and the two formulations}

\textbf{Can RealQM explain the difference between a conductor and an insulator?} That is the
question of this article: whether the framework, given only the charges present and their Coulomb
interaction, distinguishes the two --- and if not, where exactly it stops.

The question is worth asking of a new framework because the established one answers it well, and it
is worth being clear about how. Band theory takes a periodic potential, solves the one-electron
problem in it, and fills the resulting bands: a partly filled band conducts, a filled band with a gap
does not. That is a theorem plus a counting argument, it places sodium and diamond correctly, and it
does so at the size of a crystal. Density functional theory then computes the potential
self-consistently and predicts real materials at useful accuracy. What both supply rather than derive
is the starting point --- a periodic potential in the first case, an exchange--correlation functional
in the second --- and the many-electron Schr\"odinger equation from which they descend costs
exponentially in the number of electrons, so it is not itself solved at the size where the question
lives.

So the interesting question about RealQM is not whether it beats that, but whether a framework built
on Coulomb alone, with nothing fitted and polynomial cost, reaches the same distinction by its own
route. The two formulations should be put side by side before either is judged.

RealQM \cite{realqm,realqmbook} represents an $N$-electron system by $N$ non-negative densities
$\psi_i$, each carrying unit charge on its own domain $\Omega_i$, the domains non-overlapping and
their interfaces free to move to where the energy is least. The energy minimised is Coulomb's:
\begin{equation}\label{eq:E}
E[\{\psi_i\},\{\Omega_i\}]=\sum_i\int_{\Omega_i}\Bigl(\half|\grad\psi_i|^2
-\sum_a\frac{Z_a}{|x-x_a|}\psi_i^2\Bigr)dx
+\half\sum_{i\neq j}\iint\frac{\psi_i^2(x)\psi_j^2(y)}{|x-y|}\,dx\,dy,
\end{equation}
subject to $\int_{\Omega_i}\psi_i^2=1$. There is no wavefunction in $3N$ dimensions, no
antisymmetry and no self-interaction; the role played by the exclusion principle is played by the
requirement that domains do not overlap. And --- of particular relevance here --- there is no smeared
continuum anywhere in the construction: everything in it is a discrete unit charge occupying a region
of ordinary space.

Minimising \eqref{eq:E} gives stationarity conditions that take the form of Schr\"odinger's
equation --- $N$ of them in three dimensions, coupled through the Coulomb potential and through the
free boundaries --- together with a condition fixing where each domain ends. This is not the
$3N$-dimensional problem rewritten or restricted: the minimisation runs over the partition
$\{\Omega_i\}$ as well as over the densities, and a condition fixing a domain's edge has no
counterpart in a formulation that has no domains. The two share their physical content --- kinetic
energy and the Coulomb interaction, no other force and nothing fitted --- and differ in what a
solution \emph{is}: one function on $3N$ dimensions, or $N$ functions on three with free
boundaries. Neither is a special case of the other.

The contrast is often compressed to ``$3$ dimensions against $3N$'', which is the most consequential
difference but not the only one, and on its own it invites the reading that one equation has simply
been rewritten in fewer variables. The physics entering is the same; four things differ:

\begin{center}
\small
\begin{tabular}{@{}lll@{}}
\toprule
 & textbook quantum mechanics & RealQM \\
\midrule
the physics    & \multicolumn{2}{l}{Coulomb interaction $+$ kinetic energy --- the same for both} \\
\midrule
the solution   & one function on $3N$ dimensions & $N$ functions on three \\
the object     & complex antisymmetric amplitude & real non-negative density \\
               & whose square is a probability   & which \emph{is} a charge density \\
the equation   & one, linear, unitary            & $N$ coupled, nonlinear \\
the domains    & none                            & free boundaries, found by the minimisation \\
\midrule
the cost       & exponential in $N$              & polynomial in $N$ \\
its reach      & a molecule                      & a solid \\
\bottomrule
\end{tabular}
\end{center}

\noindent The last two rows are what decides whether either can be brought to the question this
article asks. The cost of the $3N$ equation grows exponentially in $N$, so its reach is a molecule
and never a crystal: at the size where conduction is a meaningful question it is uncomputable, and
only the second column is available at all.

Hartree and Kohn--Sham theory work with a drastic reduction of dimensionality, and so are no longer
\emph{ab initio}: what they compute is not settled by Coulomb's law and kinetic energy alone.

One further difference decides how the next section's argument applies here. The $3N$ equation is
\emph{linear} in $\Psi$ and its evolution unitary; \eqref{eq:E}'s stationarity conditions are
\emph{nonlinear}, since each density moves in the field of the others. That matters because the
argument of \S\ref{sec:bar} --- that a closed system rings rather than conducts, so irreversibility
must be imported --- is an argument about linear unitary evolution. Nonlinearity is where
irreversibility can arise without being inserted, which is what is claimed below for the friction a
driven partition feels.

So when conduction is said not to be derivable from Schr\"odinger's equation, what is at issue is
the linear $3N$ formulation and what can be extracted from it, not the equation itself, which
appears here too. Whether the nonlinear form escapes the same argument is not settled here for a
simpler reason: as used below it is a \emph{minimisation}, with no time in it at all, so it cannot
ring or conduct or dissipate. That is a statement about this article's method, not about the
framework's prospects.

A two-component system of ions and electrons is therefore not an idealisation the theory must be
adapted to describe. It is what the theory is already made of, in a new arrangement.

Every previous test of this construction --- ionisation energies across the periodic table, molecular
binding, nuclear binding by dual Coulomb confinement \cite{realnucleus} --- has been on \emph{bound}
matter, with an attractor always present. A plasma is the first case in which large numbers of
charges of both signs are free to arrange themselves, and it is a fair question what the framework
then gives.

\section{Which of them can answer it}

What the $3N$ equation itself can contribute is limited, and the limit is worth stating precisely,
because it is not a failure of physics but of arithmetic. Kohn established that the
conductor/insulator distinction follows from the ground state alone, with no transport theory
required; what that does is convert the question into a second one, produce the ground state. Solving
the many-electron Schr\"odinger equation for that ground state costs \textbf{exponentially} in the
number of electrons, so its reach is a molecule and never a crystal, and no growth in computing
changes an exponent. Nor does the question live where the equation can be solved: whether a response is bounded or unbounded is a statement about how it grows as the system
does --- an isolated finite sample polarises and stops whatever it is made of --- so the distinction
exists only for an extended system. The systems the exact equation can reach are not ones in which
the question can be posed, and the systems in which it can be posed are not ones it can reach. The
two domains do not overlap. A criterion whose input can never be supplied, for any system to which
the criterion applies, is not an answer held in reserve; it is a description of what an answer would
look like.

What is used at the size of a solid therefore answers a different question. Band theory supplies the
periodic potential instead of computing it and counts fillings; it places sodium and diamond
correctly and NiO wrongly, a partly filled band there belonging to an excellent insulator. Density
functional theory supplies an exchange--correlation functional. A Hubbard treatment supplies $U$ and
$t$. Each returns a verdict about a model that was put in, and the distinction that was to be
explained is among its inputs.

The second formulation is not blocked in the same way. RealQM costs \textbf{polynomially} in the
number of electrons and carries no fitted quantity, so the question can be put to it at the size at
which it is meaningful, and put from scratch. That is the whole of its claim at this point: not that
it answers correctly, but that it is the one of the two that can be asked.

The rest of the article asks it. Sections~\ref{sec:bar} and~\ref{sec:scope} fix what would count as
an answer and what kind of conductor this construction could be; the remainder puts the question to
it and reports what comes back --- which is a clean verdict on one half and an obstacle on the
other.

\paragraph{The whole of the experiment, said in advance.} It is two measurements on one chain of
seven model atoms carrying one valence electron each, and the reader is owed that before the
argument that surrounds them. \emph{First:} put the chain in a uniform field, hold the partition
fixed, and read the induced dipole --- repeated at three chain lengths, this gives $\alpha$ and how
it scales. \emph{Second:} apply no field at all; displace the whole partition rigidly by $u$ and
record the energy --- this gives the periodic curve $E(u)$, whose amplitude is the corrugation and
whose steepest slope is the field at which the pattern would begin to run. Everything else is
control: the same converged states reassembled under three boundary conditions, a range of spacings
and fillings, a mesh ladder, and the requirement that a shift by one full spacing return the same
energy. The article is long because the arguments are, not because the apparatus is.

\section{The extension, and the object it is made for}

In this article we extend RealQM, developed for atoms and molecules, to \textbf{two-component
systems}: ions and electrons, both charged, both free to move, interacting by Coulomb forces and
nothing else. A metal is such a system, and the question pursued here is the elementary one such a
system poses --- whether it \emph{conducts}, and if so what carries the charge. The same
construction covers a stellar interior and the recombination era of a Coulomb cosmology, treated as
a classical plasma.

There the one exact result is the \textbf{plasma frequency}, and it is worth saying what it is.
Displace the whole electron population rigidly by $u$ against its neutralising ion background: charge
$neu$ per unit area is exposed at the two faces, the field between them is $4\pi neu$, and the
restoring force per electron is $-4\pi ne^2u$. The population therefore oscillates at
$\omega_p^2=4\pi ne^2/m$, or $4\pi n$ in atomic units, and the construction returns it exactly.

That is a check on the Coulomb bookkeeping of a two-component system, not evidence for the quantum
structure: the derivation is electrostatics and Newton's law, and any theory with the same
electrostatics returns the same number. It is quoted here to establish that the extension to ions
and electrons is set up correctly before anything harder is asked of it. The quantities that
\emph{would} test the quantum structure --- the screening length and the dispersion of that
oscillation --- are where the construction is found wanting, in \S\ref{sec:coefficient}.

\section{What conduction asks of a theory}\label{sec:bar}

Before asking whether this construction conducts it is worth being exact about what the question
demands, because the obvious standard is one that no framework meets.

\textbf{Conduction is not derived from Schr\"odinger's equation.} The equation is unitary and
reversible while resistance is dissipative, and a closed quantum system in a perfect periodic
potential, placed in a constant field, reaches no steady current at all: $k$ runs to the zone
boundary, Bragg-reflects and returns, so the system rings --- \emph{Bloch oscillations} --- rather
than conducts. A steady current requires irreversibility, and every treatment imports it rather than
deriving it: a collision term in Boltzmann transport, a partial trace in an open-system treatment,
thermalising reservoirs in Landauer's, $\tau$ as an input in $\sigma=ne^2\tau/m$.

\textbf{What it does supply is the states.} Bloch's theorem gives extended solutions in a periodic
potential, and Pauli filling then gives the metal--insulator \emph{classification} directly: a
partly filled band conducts, a filled band with a gap does not. That is the standing of the question
below --- whether this framework possesses states of the kind that fill.

So the honest question comes in two tiers, and only the first is a fair test of a structural theory:

\begin{center}
\small
\begin{tabular}{@{}ll@{}}
\toprule
tier & status\\
\midrule
the classification: which materials conduct & from bands and filling; \emph{this is the fair bar}\\
the value of $\sigma$ & needs added statistical machinery in every framework\\
\bottomrule
\end{tabular}
\end{center}

\noindent This article asks the first question of RealQM. It does not ask the second, and no
comparison here should be read as faulting the construction for failing to supply a conductivity,
since nothing supplies one unaided. What it does examine is whether the framework can reach the
classification at all --- whether it possesses states of the kind that fill into bands --- and the
answer, developed below, is that it does not, for reasons that are structural and that the
framework's own postulate makes explicit.

\paragraph{Which clock is missing.} The first tier needs no time --- that is Kohn's point, and it is
what lets a static calculation classify. The second does, and the objection to be met is that a
minimisation could hardly report a current whatever it was applied to, copper included. So it is
worth being exact about which time is absent rather than calling the framework static. A wire
contains three time scales, twenty orders of magnitude apart, and the construction lacks them for
three different reasons:

\begin{center}
\small
\begin{tabular}{@{}lll@{}}
\toprule
what moves & its scale & why this framework does not have it\\
\midrule
the signal, hence the energy & $\sim c$ & \eqref{eq:E} is instantaneous: no retardation, no Maxwell\\
the carriers & mm/s drift & this \emph{is} $\sigma$: inserting it assumes the answer\\
the electron state & $\hbar/E_h\approx24$~as & a real $\psi$ has no phase to evolve\\
\bottomrule
\end{tabular}
\end{center}

\noindent Only the last is intrinsic --- $c$ arrives by coupling to Maxwell and drift is a transport
coefficient --- but restoring it does not by itself produce a conductivity. Attosecond variation is
far below the resolution of any conduction measurement: scattering times are femtoseconds to
picoseconds and a conductivity is an average over many of them, so the phase evolution is integrated
out before an ammeter sees anything. What it supplies is not the rate but the \emph{object} --- a
current, $J\propto\Im(\psi^*\grad\psi)$, which is identically zero for a real density and cannot be
made nonzero by any refinement of a minimisation. What is absent is a \emph{velocity} variable, and
that is a weaker requirement than a complex amplitude: in Madelung's hydrodynamic form
\cite{madelung} the same content is carried by a real density together with a real velocity field,
with $J=\rho v$ and a continuity equation in place of the phase. The vanishing of
$\Im(\psi^*\grad\psi)$ is therefore a statement about a minimisation, which has no velocity in it at
all, and not about real densities, which can perfectly well flow. The rate still comes from the
dissipation of the table above.

These two absences are one. A real non-negative density has no phase, and the same lack makes the
theory a minimisation rather than an evolution \emph{and} makes the current within a domain vanish
identically. That has a positive consequence, and it is the one this article follows: if charge
cannot flow inside a domain, the only object left that can carry it is the boundary between domains.
It also fixes how the question must be put. A framework with a clock answers ``does it conduct'' by
driving the system and reading a current; this one cannot, so the question is recast as a static one
--- the maximum gradient of $E(u)$, and the scaling of $\alpha$ with $N$ --- which is what
\S\ref{sec:threshold} and \S\ref{sec:scaling} do. The missing clock does not supply the verdict; it
determines the form in which a verdict is available at all.

\paragraph{And a pseudo-time is not a clock.} There is a time parameter in the solver, and it is
worth saying why it is not this one. The fields relax in \emph{imaginary} time, $\tau=it$, under
which $e^{-H\tau}$ damps every excited component as $e^{-(E_n-E_0)\tau}$: a filter that finds the
ground state, not a dynamics that evolves a state, and the oscillation is not lost by accident but
removed on purpose. Charge conservation goes with it. In real time $\partial_t\rho+\grad\cdot J=0$
holds exactly; continued to imaginary time the same equation becomes a diffusion with a source, with
no flux to conserve and the norm restored by a renormalisation that is global. The relaxation is
therefore not an underpowered dynamics that a smaller step would repair --- it is the branch of the
equation from which the current has been removed.

What does conserve charge locally, in the scheme as it stands, is the interface rule: cells change
ownership between neighbouring domains by the local density balance
$(\rho_i-\rho_j)/(\rho_i+\rho_j)$, so what leaves one domain arrives in the next. That is a genuine
local flux, and it is the reason the boundary and not the electron is the carrier here. It is also
where a time-dependent formulation should begin --- a continuity law with the boundary velocity as
the transport variable --- rather than in promoting a descent that was never a conservation law. The
descent itself can be promoted in exactly two ways, neither of which answers the question. Overdamped, $\partial_t\Omega=-M\,\delta E/\delta\Omega$, gives a
boundary velocity proportional to the driving force --- Ohm's law with the correct structure, except
that the mobility $M$ is put in by hand and is the conductivity up to geometry. Inertial, giving the
moving interface a mass, gives oscillation: the chain rings, as any closed system does, and carries
no steady current. The descent rate is a numerical parameter with no calibration to anything
physical, and each way of calibrating it either inserts the answer or returns the ringing.

\section{What kind of conductor this can be}\label{sec:scope}

The construction does not have a free choice of target, and identifying it correctly settles most
of what follows. Non-overlap fixes two things at once, and they happen to be the two parameters of the Hubbard
Hamiltonian. The on-site repulsion $U$ --- what it costs to put two electrons on one site --- is
infinite, since two electrons may not share a domain at all; and the transfer integral $t$ --- the
amplitude to hop to a neighbour --- is exactly zero, since it \emph{is} the overlap of neighbouring
orbitals and there is none. In that
language the construction is the \textbf{atomic limit} and not the Mott limit --- the reading it is
usually given --- and the difference is not a refinement:

\begin{center}
\small
\begin{tabular}{@{}lll@{}}
\toprule
 & $U\to\infty$, $t\neq0$ & $U\to\infty$, $t=0$\\
\midrule
half filling      & Mott insulator      & isolated atoms\\
superexchange     & $J=4t^2/U$          & none\\
doped             & \textbf{conducts}   & still nothing\\
spin dynamics     & yes                 & none\\
\bottomrule
\end{tabular}
\end{center}

\noindent A Mott insulator is insulating but alive: it carries magnetic exchange, and doping it
produces a conductor, which is what the cuprates are. The atomic limit is a collection of atoms with
electrostatics between them. This accounts for every negative result below, and for one in
particular: doping a Mott insulator works \emph{through} $t$, so at $t=0$ there is nothing for
doping to activate --- which is exactly what the measurement shows, the barrier rising with carrier
concentration rather than falling. Metals are therefore out of scope in kind and not by degree, their carriers being extended states
and their transport band motion, neither of which an arrangement of densities on domains appears
able to supply. Two results usually read as failures follow from that scope rather than against it: $d\sigma/dT>0$ is \emph{correct} for an
activated conductor and fatal only for a metal, and the absence of a Fermi surface does not bear on
a dilute carrier gas, which is non-degenerate and has none to reproduce.

What remains within reach is matter whose carrier is a localised charge on a site, whose gap is
on-site repulsion rather than band structure, and whose transport is activated hopping. That is a
real class --- doped Mott insulators, organic and molecular semiconductors, polaronic oxides,
amorphous semiconductors --- and this article set out to claim it. \S\ref{sec:conduction} measures
the activation energy and finds $613$--$1723$~meV, against the $10$--$500$~meV those materials
show. The claim does not survive its own measurement, and what the construction describes is the
insulating end of the class and not the conducting one. Silicon and the III--Vs were never in
scope, their gaps being band gaps.

\section{The interface, and the one number that is wrong}\label{sec:coefficient}

The free interface has a defect of its own, recorded here because it conditions everything the
partition does, though the hopping barrier does not turn on it. Partitioning a gas of density $n$
gives a kinetic energy density $t(n)=C\,n^{5/3}$ --- the Thomas--Fermi \emph{form}, obtained from
partition geometry rather than Fermi statistics --- but the flat trial function is admissible on a
free boundary, so $C=0$ exactly. There is no degeneracy pressure, the screening length vanishes
rather than being finite, and the plasmon does not disperse. The defect is one number and not the
construction: $C$ is fixed by the interface condition, and a Robin condition at $\beta a=1.146$
gives $C=2.871=C_{\rm TF}$ exactly, against Dirichlet's $14.804$, all three carrying the correct
$n^{5/3}$ scaling.

Whether that value is tolerable for matter is the open question, and it is sharper than a parameter
check. Adding the surface term and scanning it, $\mathrm{H}_2$ unbinds by $\beta a\approx0.71$,
below the $1.146$ the electron gas requires: Teller's theorem, that Thomas--Fermi binds no molecule,
is met inside the framework rather than inherited from outside it. Double occupancy --- the repair
conduction would need at half filling --- becomes favourable at $\beta a=0.451$ and is affordable by
comparison. These figures were obtained at differing densities and are not strictly commensurable;
the equation of state is pursued elsewhere, and what matters here is only that the interface
carries no cost of its own unless one is supplied.

\paragraph{A caution about specified partitions, learned by getting it wrong.} The method of
\S\ref{sec:conduction} evaluates the energy of a \emph{stated} partition, and it is sound only when
the stated partition is one the physics admits. The interface is not a surface to be shaped at
will: the wavefunction is a single global field partitioned by labels, continuous across every
interface, each domain carrying one unit of charge, and the boundary sits where the densities
balance. Freezing a partition removes that last condition. An attempt to measure a bending rigidity
by imposing a curvature, freezing it and relaxing the fields gave an energy falling symmetrically
and quadratically in the imposed curvature, which reads convincingly as a negative rigidity; but
released, a boundary started from the curved configuration returns to the flat one. The flat
interface is the balance surface and is stable, and what the frozen scan measured was the energy
gained by \emph{dropping the balance condition}. A specified partition gives a legitimate energy
difference between configurations the physics admits --- as symmetry guarantees for the two in
\S\ref{sec:conduction} --- and a meaningless one along a direction it forbids. Nothing in the
numbers distinguishes the two cases; the check is to release the boundary and see whether it stays.

\section{Computing with it anyway}\label{sec:conduction}

Electric conduction is the interaction of the electrons with the ion lattice. It is the elementary
case --- ordinary matter at ordinary temperature --- and it is for that reason the \emph{demanding}
one here, since ordinary-temperature conduction is degenerate matter, and degeneracy is what this
framework has least of.

One property of the construction is needed before anything else, and it is the property that makes
conduction a live question rather than a hopeless one. With the free interface used throughout, a
uniform density is admissible on any partition: $\psi_i\equiv|\Omega_i|^{-1/2}$ satisfies the
normalisation, minimises the electrostatic energy against a neutralising background, and has
$\grad\psi_i=0$. Redrawing the boundary therefore costs nothing,
\begin{equation}\label{eq:Czero}
\text{a free interface between domains of equal density carries no confinement energy,}
\end{equation}
and the boundary is the only object this framework moves --- its solver evolves each interface by the
local density balance $(\rho_i-\rho_j)/(\rho_i+\rho_j)$, transferring ownership of a cell when the
neighbour's density exceeds its own. Domains do not translate; they are redrawn. \emph{A vanishing
interface cost is a vanishing cost of transport.} What that same fact does to the equation of state
is a separate matter and is taken up in \S\ref{sec:coefficient}, where it is much less welcome.

Metallic conduction is the paradigm of \emph{delocalisation}: the carriers are extended states
spanning the crystal, whereas RealQM's electrons are localised by construction, each in its own
domain, with non-overlap requiring $N$ carriers to occupy $N$ disjoint regions. The natural minimiser
of \eqref{eq:E} is a set of compact cells, which describes a Mott insulator more naturally than a
conductor. Nor is there an evident mechanism for the metal--insulator distinction, which in standard
theory is band filling --- two electrons per band per cell --- an exclusion-principle count the
framework does not have; and with no Fermi surface, the carrier density in $\sigma=ne^2\tau/m$ has no
referent.

What does fit naturally is the opposite regime: localised charges hopping between domains by
thermally activated jumps, as in semiconductors, disordered systems and small-polaron transport. That
needs no Fermi surface and follows from the same thermal motion that supplies the pressure.

This yields a prediction crude enough to be decisive, and it concerns only a sign. Metallic
conductivity \emph{falls} with temperature; hopping conductivity \emph{rises}. If RealQM transports
charge by hopping rather than band motion it should give $d\sigma/dT>0$. Within the scope set in
\S\ref{sec:scope} that is the \emph{right} answer, activated transport being what a localised-carrier
semiconductor does; it would be fatal only for a metal, which is why the target matters. Whether it does has not been computed for electrons, and the
pessimistic reading should be held loosely: the same lattice that reproduces alkali lattice constants is
present here too.

For \emph{protons} the mechanism has already been demonstrated in this framework. Grotthuss
conduction --- an excess proton migrating along a hydrogen-bonded chain of water molecules by
successive transfers, rather than by bodily motion of any one particle --- is hopping transport in
its purest form, and it is a standard RealQM calculation: a chain of waters with $\mathrm{O}$--$\mathrm{O}$
spacing near $2.65$~\AA{}, an excess proton, and the question whether it hops from one molecule to
the next. Related transfers, $\mathrm{HF}+\mathrm{H_2O}\to\mathrm{F^-}+\mathrm{H_3O^+}$ and
$\mathrm{HCl}+\mathrm{NH_3}\to\mathrm{Cl^-}+\mathrm{NH_4^+}$, are computed the same way.

This matters for the argument rather than merely illustrating it. Charge transport by transfer
between localised domains is not a regime the framework must be stretched to reach --- it is one it
already handles, in the case where the carrier is a proton.

The electron case cannot presently be tested in the same way, and the obstruction is instructive
enough to be worth setting out. Two chain calculations were attempted here --- five hydrogens with an
excess electron, and five hydrogen kernels with one electron missing, the electronic counterparts of
the proton wire --- and both returned nothing. Neither result means what it appears to.

The cause is not what it first appeared. Kernel proximity assigns the labels only \emph{initially};
thereafter a free-boundary rule evolves each interface from the local density balance and transfers
ownership of a cell when the neighbouring domain's density exceeds its own. Partitions therefore do
rearrange, and electron transfer is representable in principle. What defeated the attempts here was
specific to each: a carrier given no kernel feels no force from the lattice; a bare kernel is given no
domain, so a vacancy has nothing to migrate into; and a carrier placed on an occupied site can seize
the nuclear region outright, since with free interfaces no confinement cost opposes it. Each is a
property of the construction, not of \eqref{eq:E}.

The deeper cause survives that repair. Conduction is charge in motion, and \eqref{eq:E} is an energy
minimisation: it returns the equilibrium partition, not a current. Given complete freedom in the
labels one would still obtain the ground state, and for a symmetric chain that state places one
domain per site by symmetry, with nothing flowing. A minimisation has no time in it.

\paragraph{Which filling the measurements belong to, and which needs the repair.} It is worth being
exact about this, because it determines how much of the article depends on $\beta$. At \emph{half
filling} --- one electron per site, every domain occupied --- the $U\to\infty$ construction forbids
conduction outright: there is no vacancy for a carrier to enter, and the insulator is structural
rather than dynamical. That is the case the double-occupancy repair is for. But it is not the case
measured here. The chains reported above are doped by ionising every second atom, leaving half the
sites empty, and every barrier in this article is the cost of a carrier moving into a \emph{vacancy}.
Single occupancy handles that without amendment.

So the transport results stand on their own: they require no finite $U$, no double occupancy and no
$\beta$. What requires the repair is the undoped, half-filled chain, which is not run here. This is
the stronger position of the two --- the measurements are independent of the proposal made to extend
them --- and it locates precisely what the proposal would buy: not better hopping, but conduction at
a filling where the construction currently permits none.

What the framework does possess is an asymmetry of dynamics. Nuclei move, under classical equations
of motion; electrons relax, to a static minimum. Charge carried by \emph{nuclei} is therefore
representable --- which is precisely why the Grotthuss calculation works, the carrier there being a
proton --- and charge carried by electrons is not, because in this scheme electrons have no equation
of motion at all. The correct machinery exists elsewhere in the programme, in the time-dependent form
where a real density carries a complex current, and a conduction calculation belongs there rather
than here.

\section{The threshold field, measured}\label{sec:threshold}

The question this article asks has a form that avoids barriers altogether, and it is the form in
which the answer can actually be obtained.

Put the chain in a field. Every electron shifts to the right. If the shift reaches one full lattice
spacing, each electron has taken over its neighbour's kernel --- and since the electrons are
identical and the energy depends on the partition and not on labels, \emph{the configuration is
restored}. That restoration is exactly what separates a current from a polarisation:
\begin{center}
\small
\begin{tabular}{@{}ll@{}}
\toprule
tilt below the steepest slope of $E(u)$ & $u$ saturates: bounded polarisation, no DC current\\
tilt above it & $u$ runs: the state cycles, charge flows steadily\\
\bottomrule
\end{tabular}
\end{center}

\noindent This \emph{is} the applied-field, free-boundary case, and it is worth saying so plainly:
$u$ is the collective coordinate of the free boundaries themselves, the displacement a field would
drive, and a uniform field adds nothing to the energy but a term linear in it, $E(u)-NeFu$. The
field therefore tilts the landscape and changes nothing else in it, so the landscape can be computed
once, without a field, and tilted afterwards. What is given up by doing it this way is the motion
\emph{after} the last stationary point is lost --- the rate, which the paragraph below shows would
be an artifact in any case --- and not the threshold, which is exactly where that point is lost.

So the threshold is not an activation energy and not a saddle. It is the \textbf{maximum
gradient} of the periodic energy $E(u)$, and \S\ref{sec:mesh} has measured that curve on the wrapped
chain, where $E(u{+}a)=E(u)$ holds to every digit and the three reflection pairs agree to
$1$--$2\times10^{-5}$~Ha. Nothing here differences distant configurations.

From the measured curve at $r_c=1.93$, $a=5.5$~a$_0$ --- both parameters determined rather than
chosen --- the steepest segment gives $|dE/du|=0.0982$~Ha/a$_0$ for seven electrons, and a
sinusoidal fit, which recovers the peak a chord underestimates, gives $0.1194$. Per electron:
\begin{equation}\label{eq:Fdep}
F_{\rm dep}=0.014\text{--}0.017~\text{au}
 =7\text{--}9\times10^{9}~\mathrm{V/m}
 =2.1\text{--}2.6~\mathrm{V\ per\ lattice\ site}.
\end{equation}

Copper carrying an ampere runs at $4\times10^{-12}$~V per site. The separation is \textbf{twelve
orders of magnitude}, and it is not a number any refinement moves. Nor is such a field available:
dielectric breakdown occurs two to three orders of magnitude below $10^{10}$~V/m, so the lattice is
destroyed long before the electron distribution depins. \emph{The construction does not conduct
under any field a material survives.}

\paragraph{Above the threshold, and why the rate there is not a result.} Drive harder than
\eqref{eq:Fdep} and the partition does run: the tilt exceeds every slope of $E(u)$, no configuration
is stationary, and the solver will advance the interfaces without stopping. It will also report a
speed, and that speed is worth naming as an artifact before anyone quotes it. What sets it is the
descent parameter of the relaxation scheme, not the physics --- as \S\ref{sec:bar} sets out, the
pseudo-time carries no calibration --- so the same chain under the same field yields whatever current
the numerical mobility is set to. The threshold is a property of the measured curve $E(u)$ and
survives refinement; the current above it is a property of the scheme and does not --- and the
regime is in any case the one just excluded.

\paragraph{Why this supersedes the barrier attempts.} Every earlier route asked for a difference
between two configurations far apart in the energy landscape --- and \S\ref{sec:mesh} shows why that
could not work, six of eight scans failing their own control. Equation \eqref{eq:Fdep} asks instead
for the slope of a single curve, sampled at points a fraction of a spacing apart, on a geometry
whose control passes by four orders of magnitude. It is the same physics, asked in the form the
discretisation can answer.

\paragraph{A defect in the model behind these numbers, and what it leaves standing.} The
pseudopotential exclusion used throughout this section keeps valence density out of a domain's
\emph{own} core only. A cell belonging to domain $A$ but lying inside neighbouring core $B$ is not
excluded, so $A$'s density enters $B$'s core and collects its attraction. That is invisible wherever
each domain contains its own core --- which is why $d=0$ is unaffected --- and decisive at
$d=\tfrac12$, where every core is split between two domains and the exclusion never fires on either
half. Excluding from \emph{every} core instead inverts the registry: the minimum moves to $d=0$,
atom-centred, and the bond-centred preference reported here disappears.

Neither treatment is usable as it stands. Own-core exclusion under-excludes and gives the wrong
registry; full hard exclusion over-confines, since at $r_c=1.93$ and $a=5.5$ the excluded spheres
leave only $1.64\,a_0$ between neighbours and the electron at $d=\tfrac12$ becomes a particle in a
box rather than a charge in a lattice. The two therefore \emph{bracket} the answer rather than
settling it, and what is needed is a soft repulsive core, or $r_c$ recalibrated with full exclusion
active. \textbf{$V_0=816\pm4$~meV and the threshold field \eqref{eq:Fdep} are the own-core values},
verified as such but resting on an exclusion rule that is not the model. Refinement cannot repair
this: the mesh convergence below is convergence of the wrong quantity.

The scaling route of \S\ref{sec:scaling} does not inherit the defect. Its measurements are made at
an undisplaced partition, where each domain contains its own core and the two rules agree to five
decimals, so $\alpha=27.4N$ and the extensivity verdict stand independently of how the exclusion is
settled. That the article reaches the same conclusion by two routes is what makes the loss of one of
them survivable.

\paragraph{And it is converged.} The amplitude was checked on three grids at fixed geometry,
$a=5.5$, $r_c=1.93$, with the translation control exact and the reflection pairs identical to five
decimals at each:

\begin{center}
\small
\begin{tabular}{@{}lrrr@{}}
\toprule
$h$ (a$_0$) & $0.448$ & $0.352$ & $0.299$\\
$V_0$ per electron (meV) & $812$ & $820$ & $815$\\
\bottomrule
\end{tabular}
\end{center}

\noindent $V_0=816\pm4$~meV, a spread of half a per cent across a fifty per cent refinement. What
makes that possible is worth stating, because the same solver produced scatter of a factor of two
in \S\ref{sec:mesh}: over this range the \emph{totals} drift by $1572$~meV while the amplitude moves
by $7$. The cancellation is $99.5\%$, against roughly $89\%$ for the open-chain barrier --- and
$89\%$ left a residual comparable to the signal, where $99.5\%$ leaves one per cent of it. The
improvement comes from the geometry alone: $d=0$ and $d=\tfrac12$ lie on one periodic curve on an
identical grid with identical topology, so their discretisation errors track each other, whereas the
open chain's two configurations differed in their end domains as well as their registry and the
errors were independent. \emph{A difference is only as good as the symmetry relating the two
configurations}, which is the transferable lesson of this article.

\paragraph{The corrugation is bounded below, at every density.} The same scan repeated across
lattice spacings, with the translation control exact ($0.0$~meV) at every one and reflection pairs
agreeing to $2$--$3\times10^{-5}$~Ha:

\begin{center}
\small
\begin{tabular}{@{}lrrrrr@{}}
\toprule
$a$ (a$_0$) & $3.5$ & $4.5$ & $5.5$ & $6.5$ & $8.0$\\
$V_0$ per electron (meV) & $1513$ & $972$ & $812$ & $\mathbf{783}$ & $816$\\
\bottomrule
\end{tabular}
\end{center}

\noindent $V_0$ has a \emph{minimum} near the equilibrium spacing and rises on both sides:
compressed, the cores crowd the electron; expanded, sharing a kernel at the bond-centred registry
costs more than owning one. So the corrugation is bounded below at ${\approx}783$~meV per electron
across a factor of $2.3$ in lattice constant, never approaches zero, and the registry preference
never flips --- there is \textbf{no free-sliding point at any density}.

That settles a question the structural argument leaves open. A framework that could reach $V_0\to0$
would produce both phases and could decide conduction on its own terms; this one produces only the
pinned phase. And the trend runs the wrong way: real matter \emph{metallises} under pressure, as
screening flattens the corrugation, whereas here compression makes the pinning worse. This also
restores the compression result reported earlier by the free-boundary route which
\S\ref{sec:mesh} discredits --- a ratio between configurations sharing partition family and mesh is
far more robust than $V_0$'s absolute value.

\paragraph{What is bounded, and what is not.} Everything measured here is the \emph{rigid} slide:
all $N$ electrons displaced together, at half filling, where the conducting electrons are also the
binding electrons. That is the expensive path, and its cost grows with chain length. The cheap path
is a single carrier moving while the rest hold still --- a kink, whose cost is length-independent
--- and that is the configuration real conductors use, a supporting lattice that binds and a dilute
carrier that conducts. \emph{The carrier barrier is not measured in this article.} A kink on the
cyclic geometry requires the slip to accumulate to exactly one unit around the loop,
$f(j+N)=f(j)+1$, which the profile used here does not satisfy; the attempt fails its own control by
$997$~meV and is not reported as a result. Note also that the maximally delocalised limit of such a
profile gives evenly spaced walls and a barrier of exactly zero, so the kink barrier is a function
of its localisation and the physical value is the minimum over that --- a two-parameter
determination, not a single number. \textbf{The figures above are therefore upper bounds on the
quantity conduction depends on}, and the verdict rests on them together with the structural
argument, not on a measured carrier barrier.

\paragraph{A third qualification, and the largest.} Every energy in this article is computed
with the \textbf{nuclei clamped} at their lattice sites. So $V_0=812$~meV is the rigid-lattice
corrugation, and \eqref{eq:Fdep} the rigid-lattice threshold --- both upper bounds, because in real
matter the lattice relaxes around the carrier and that relaxation is the dominant term. It is what
makes small polarons work at all: the Marcus--Holstein barrier is $E_a=\lambda/4$ with $\lambda$ the
reorganisation energy, and $\lambda\approx0.1$--$0.3$~eV in organic semiconductors gives
$25$--$75$~meV, far below the bare electronic corrugation of those same materials. Nothing here
tests it, and the framework is not prevented from doing so --- nuclei move in this scheme under
classical equations of motion, which is the half of the dynamics it possesses and the reason the
Grotthuss calculation works. If lattice relaxation reduced the corrugation by the factor it does in
polaronic materials, $812$~meV would fall toward $100$~meV and the threshold with it, from
${\sim}8$ to ${\sim}1$~GV/m: still above dielectric breakdown, but no longer by three orders. This
is the one accessible variation that could move the verdict rather than confirm it, and it is
untested.

The two bounds are independent and they multiply, which is worth stating in one place rather than
leaving to be assembled: a clamped lattice and a rigid slide both raise the number reported, the
first by an estimated order of magnitude and the second by a factor nobody here can quote, since the
kink barrier is the quantity \S\ref{sec:mesh} shows cannot be extracted. Twelve orders is therefore
the separation for the mode measured, not a floor under every mode. What survives either way is the
sign of the conclusion --- the field required is far above the one at which a wire carries current,
and the reduction would have to reach nine orders to change it.

\paragraph{Two qualifications, both real.} This is \emph{rigid} sliding, every electron moving at
once, which is the expensive path --- the electron spacing is held fixed, the whole pattern
translating without deforming --- so a kink, which compresses the domains ahead of it and expands
those behind, depins at a lower field and \eqref{eq:Fdep} is an upper bound. And the chain is unstable to non-uniform distortion --- a sinusoidal displacement
\emph{lowers} the energy, giving an effective stiffness $\kappa<0$, because domains here cannot
change position without changing width and a wider domain lets its electron spread. What the
distorted state does under a field is not determined here. Both qualifications point the same way,
toward a lower threshold, and neither is remotely capable of closing twelve orders of magnitude.

\paragraph{Why a chain rather than a ring.} A ring supplies a clean barrier --- no ends,
every site equivalent, the energy obliged to repeat with the lattice period --- but it cannot host
the drive: a conservative field has zero circulation, so no bias drives a sea around a closed loop.
A chain has ends, a potential drop along it is an ordinary thing to apply, and a single carrier may
enter at one end and leave at the other, which is the transport event itself.

\paragraph{The partition must be specified, not searched for.} Before the number, the method,
because the obvious approach fails and the failure is instructive. A kink is \emph{not} the ground
state, so a free-boundary minimisation relaxes it away --- asked to hold one, the solver oscillates
between the kinked and uniform partitions and the energy never settles, wandering over
$0.05$~Ha against a barrier of order $0.01$. That is not a convergence failure to be tuned away:
it is unaltered by the surface tension $\beta$, by reducing the boundary step fivefold, and by
freezing the labels outright.

The cause is general and it disqualifies the free boundary for every transport quantity in this
article. The interface moves by \emph{local single-cell label flips}, so a run descends to whichever
configuration no single flip improves --- the nearest stuck point, not the minimum --- and which one
that is depends on the initial condition. Four independent geometries agree:

\begin{center}
\small
\begin{tabular}{@{}lll@{}}
\toprule
geometry & spread & against a barrier of\\
\midrule
ring, between two builds & $0.005$~Ha & $0.012$--$0.022$\\
kink chain, across steps & $0.05$~Ha (never settles) & $\sim0.01$\\
$\mathrm{H}^-+\mathrm{H}$, four seeds & $0.031$~Ha & $\sim0.001$\\
alkali chain, two seeds & $0.013$~Ha & $\sim0.001$\\
\bottomrule
\end{tabular}
\end{center}

\noindent In every case the ambiguity exceeds the quantity sought, by one to two orders. It is a
failure of the \emph{search}, not of the energy functional, which is why nothing in the functional
touches it. It is also why the ring's $0.01186$~Ha moved to $0.022$~Ha under a recompilation that
altered no physics: rounding at the last bit flips cells sitting near the threshold, and a different
stuck point follows. The rigid-sliding figure should be read as a bound of that order rather than a
determination, and the doping barriers of the earlier barrier scans, obtained the same way, with it.

The repair is not a better minimiser but no minimiser. The domain walls of a chain are planes
perpendicular to the axis, so they can be \emph{placed}: specify the partition, freeze it, and relax
only the wavefunctions. Nothing is searched for, the same input returns the same output, and the
question asked is no longer ``where does the boundary go'' but ``what does this configuration
cost'' --- which is the question a barrier actually poses. It is the device that makes the
$\mathrm{H}_2$ midplane exact, applied to a coordinate instead of a symmetry.


\section{The same verdict from scaling}\label{sec:scaling}

The measurements in this section are on a \emph{different chain} from \S\ref{sec:threshold}, and
that should be said before any of its numbers appear. It is a model atom of nuclear charge $+1$
with core radius $r_c=1.0\,a_0$ and one valence electron, seven of them in a row at spacings from
$3.5$ to $8.0\,a_0$ --- parameters chosen to define a system rather than to represent an element.
The row is an arrangement of kernels, not a property of the construction: the domains are regions of
ordinary three-dimensional space, and the same computation runs unchanged on a square or cubic
lattice of the same atoms. Its corrugation runs from
$182$ to $461$~meV per electron across that range, and is $379$~meV at $a=4.5\,a_0$, the spacing at
which the polarisabilities below are computed. Nothing here is a second determination of the
quantities of \S\ref{sec:threshold}; what is claimed is a statement about \emph{scaling}, which is
independent of which chain it is made on.

That independence is worth making concrete, since the natural objection is that the verdict rests on
two chains rather than one. Extensivity does not involve $r_c$ at all: \eqref{eq:extensive} is
proportional to $N$ for \emph{any} positive stiffness, so $r_c$ can move $\alpha_{\text{intra}}$ and
$K$ --- both per-atom constants --- without touching the power of $N$. What $r_c$ could do is take
the one escape, $K\to0$, and the chain of \S\ref{sec:threshold} is measured to be further from it
than this one: $V_0=812$ against $379$~meV per electron, hence $K\simeq530$ against
$369$~meV/a$_0^2$ and a slide term $e^2/K\simeq51$ against $73.6\,a_0^3$. The stiffer chain is the
safer one. Recomputing $\alpha$ at $r_c=1.93$ would therefore move the coefficient $27.4$, which is
reported as a measurement of this chain and not as a property of the framework, and would leave the
classification where it stands.

The threshold field settles the question by driving the pattern. It can also be settled without
driving anything, by asking how the response grows with the size of the sample --- an independent
route, resting on different assumptions, and worth having because the two can be checked against
each other.

One point about what is being asked deserves stating before any of it, because it decides which
part of the measurement carries the weight. An isolated finite system cannot conduct, whatever it is
made of. Place a piece of copper in a box and apply a static field: the charge redistributes to the
surfaces, the interior field cancels, and everything stops. Copper polarises and stops. So the fact
that a response \emph{saturates} is never evidence of insulation --- it is what any closed system
does, and no calculation on a bounded sample can conclude otherwise.

What distinguishes the two cases is not whether the response saturates but how the saturated value
\emph{scales}. A metal's polarisability grows faster than the number of atoms; an insulator's is
extensive. That is what Kohn's argument supplies and what makes a closed, static calculation
capable of answering the question at all. The measurement that carries the classification here is
therefore $\alpha = 27.4N$ --- the extensivity, not the boundedness --- and the corrugation enters
only because a vanishing $K$ is the one way extensivity could fail.

\subsection{Why the scaling is fixed before it is measured}

Nor is the outcome fixed in advance by the postulate. At full filling the domains tile space, but
transport does not require a vacancy: every boundary may shift at once, translating the whole
partition. That collective mode has a polarisability. Writing $K$ for the per-electron curvature
of the corrugation the partition slides through,
\begin{equation}
E(u) = \tfrac12 K_{\text{tot}}u^2 - NeFu
\quad\Longrightarrow\quad
\alpha_{\text{slide}} = \frac{N^2e^2}{K_{\text{tot}}} = \frac{Ne^2}{K},
\label{eq:slide}
\end{equation}
since the stiffness is extensive, $K_{\text{tot}}=NK$.

Each domain carries exactly one electron. The constraint $\int_{\Omega_i}\psi_i^2 = 1$ is imposed,
not obtained, so the monopole of every domain is pinned and \textbf{no charge passes between
sites}. The sites, however, are not the fixed points of the problem --- the kernels are. A domain
may arrive at the kernel that was its neighbour's while carrying the same unit of charge it began
with, so transport here is \emph{the sites moving}, not charge moving between them, which is the
dislocation relation again and the reason nothing needs to hop. That removes the mechanism which gives extended systems their large polarisability, and
leaves exactly two responses to a field: the density deforming within each fixed domain, and the
domain boundaries displacing. Write $\alpha_{\text{intra}}$ for the first --- the polarisability an
atom would have with its domain held rigid, the electron leaning toward the field while its walls
stay put --- and use \eqref{eq:slide} for the second:
\begin{equation}
\alpha = N\Big(\alpha_{\text{intra}} + \frac{e^2}{K}\Big),
\label{eq:extensive}
\end{equation}
which is proportional to $N$ exactly. Allowing the walls to move differentially does not change
this: with elastic coupling $\kappa$ between walls and substrate stiffness $K$ per site, the
displacement profile satisfies $\kappa u'' - K u = -f$ and returns $u \simeq f/K$ through the bulk
with corrections over a length $\sqrt{\kappa/K}$ at the ends, so the total dipole remains
proportional to $N$.

The one escape is $K \to 0$. Only when the corrugation vanishes is the displacement limited by the
chain length rather than by a restoring force, and only then does $\alpha$ grow faster than $N$.

\subsection{Polarisability against chain length}

Computed for a model chain --- one-valence-electron atoms of core
radius $r_c=1.0\,a_0$ at spacing $a=4.5\,a_0$, frozen partition, from the induced dipole at two
field strengths well below the depinning threshold:

\begin{table}[h]
\centering
\begin{tabular}{rrrr}
\toprule
$N$ & $\alpha$ & per atom & linearity \\
\midrule
3 & 81.57 & 27.19 & 4.3\% \\
5 & 137.26 & 27.45 & 5.5\% \\
7 & 192.50 & 27.50 & 4.8\% \\
\bottomrule
\end{tabular}
\caption{Polarisability against chain length. The last column is the spread between the two field
strengths. It is the precision of the measurement rather than a departure from linearity: the
fields are already at a tenth of the depinning threshold, and lowering them further puts the
induced energy below the solver's own convergence noise, so $5\%$ is the floor available here.}
\label{tab:chainN}
\end{table}

The polarisability per atom is constant to $1\%$ across the three lengths, so
\begin{equation}
\alpha = 27.4\,N,
\label{eq:alphafit}
\end{equation}
reproducing all three entries to better than $1\%$ with no intercept required. The response is
extensive, with $\alpha_{\text{bulk}} = 27.4\,a_0^3$ per atom.

A surface term is not resolved. At the $5\%$ precision of these values an intercept anywhere
between roughly $\pm 10$ is admissible, so the ends of the chain cannot be shown to differ from its
interior --- which is consistent with the derivation, since a constraint that fixes the charge on
every domain makes no distinction between an end site and a middle one, but it should be read as
an absence of evidence rather than evidence of absence.

This is \eqref{eq:extensive} with its coefficient supplied, and it confirms rather than tests the
derivation: extensivity follows from the postulate, and what the measurement adds is the size of
the coefficient and the fact that nothing anomalous happens as the chain lengthens.

One caution belongs with that coefficient, and it is not small. The partition is frozen, so
$e^2/K$ drops out of \eqref{eq:extensive} entirely and what is measured is $\alpha_{\text{intra}}$
alone. Taking $K$ from the corrugation at the same spacing gives $e^2/K = 73.6\,a_0^3$,
nearly three times $\alpha_{\text{bulk}}$. Releasing the partition would therefore roughly triple
the response: the wall displacement is the \emph{larger} part of it, and the numbers above are a
lower bound on the polarisability rather than an estimate of it.

It was frozen for a reason of resolution and not of convenience. The solver moves an interface by
transferring ownership of whole cells, so a wall sits at a cell boundary, while a sub-threshold field
asks it to move by a fraction of one: it stays put, or jumps the full cell. A released measurement
returns a staircase rather than a linear response. That is also why \S\ref{sec:threshold} imposes
$u$ rigidly instead of letting a field find it --- an imposed displacement can be set to any value
and checked exactly, since translation by one spacing must return the same energy. Releasing the
partition and re-measuring $\alpha$ directly, which would test the factor of three above, needs a
sub-cell representation of the interface and has not been done.

\subsection{What the boundary treatment fixes, and what it does not}

$V_0$ is a difference between two \emph{wall placements}, so of every quantity reported here it is
the one most exposed to how boundaries enter the kinetic energy. That exposure was measured rather
than assumed. Reassembling the same converged states with the surface condition changed --- the
default Neumann at walls and cores, then Dirichlet at the walls, then Dirichlet at the cores ---
gives $+379$, $+721$ and $-949$~meV per electron at $a=4.5\,a_0$.

Two readings follow, and they point opposite ways. The corrugation is \emph{surface-dominated}: the
bulk kinetic energy is nearly unchanged while $V_0$ moves by $1300$~meV and reverses, which is
direct support for the steric character argued below, since a quantity fixed at the exclusion
surfaces is by definition not an electronic one. But it also means neither the magnitude nor the
sign of $V_0$ can be read as a property of the model alone. The Dirichlet variants are a stress
test rather than a rival physics --- forcing the density to vanish across one cell adds a large and
mesh-dependent surface term --- yet the sensitivity they expose is real.

What survives is what the classification needs. $|V_0|$ is $379$, $721$ and $949$~meV in the three
treatments; it approaches zero in none of them, so $K$ stays finite and the polarisation bounded
however the surfaces are handled. The verdict is robust where the numbers are not.

\subsection{Two routes, one verdict}

These are different chains, as stated at the head of the section, so the agreement claimed is
qualitative --- that both give a bounded response and a positive curvature --- and not a numerical
coincidence between two determinations of one quantity.

The threshold field of \S\ref{sec:threshold} and the extensivity measured here are independent. One
drives the partition and reports the field at which it would run; the other never applies a field
large enough to move anything and reports how the bounded response scales. They rest on different
assumptions --- the first on a periodic energy curve and its steepest gradient, the second on the
Kohn--Resta criterion and the growth of $\alpha$ with $N$ --- and they agree: the polarisation is
bounded, the curvature is positive, and there is no DC current. That agreement is the strongest
statement this article makes about the classification, because neither measurement could have
produced it alone.


\section{The barrier is not extractable, and the control that shows it}\label{sec:mesh}

A barrier taken as the difference of two configurations is only as good as the cancellation between
their errors, and that cancellation is not something one may assume. The check is available and
costs nothing conceptually: scan the carrier's position through a \emph{full lattice period} rather
than sampling two configurations, and require the scan to return to where it began. On an infinite
chain $E(c{=}0)$ and $E(c{=}1)$ are the same configuration translated by one spacing and must agree
exactly. Their mismatch is therefore a per-run noise floor --- finite size and discretisation
together --- measured in the same run as the quantity it has to be compared against.

Two further benefits come with it. The barrier becomes the \emph{amplitude} of $E(c)$, so nothing
depends on having guessed where the extrema lie; and the shape of $E(c)$ is visible, so a genuine
corrugation can be told from a monotone drift. Both matter, as it turns out: the maximum is
generally not at the half-way point that a two-configuration measurement assumes.

Applying it is decisive. Eight scans, each nine configurations over one period, with the core
radius and the lattice spacing varied:

\begin{center}
\small
\begin{tabular}{@{}rrrrrl@{}}
\toprule
$a$ (a$_0$) & $r_c$ & amplitude & control & ratio & \\
\midrule
$6.0$ & $1.60$ & $\phantom{0}770$~meV & $699$~meV & $\phantom{0}1.1$ & fails\\
$6.0$ & $1.60$ & $1087$ & $535$ & $\phantom{0}2.0$ & fails\\
$4.5$ & $0.80$ & $1854$ & $1524$ & $\phantom{0}1.2$ & fails\\
$4.5$ & $1.20$ & $1735$ & $936$ & $\phantom{0}1.9$ & fails\\
$4.5$ & $1.60$ & $1900$ & $\phantom{00}78$ & $24.4$ & passes\\
$4.5$ & $1.93$ & $2597$ & $584$ & $\phantom{0}4.4$ & passes\\
$5.5$ & $1.93$ & $1949$ & $1192$ & $\phantom{0}1.6$ & fails\\
\bottomrule
\end{tabular}
\end{center}

\noindent \textbf{Six of the eight fail their own control}, and the two that pass do not survive a
change of either parameter. Worse for any attempt to rescue a number from them, the passes cluster
at one spacing, $a=4.5$, across three different core radii --- including one for which $4.5$ is not
the equilibrium. Had the control been passing because a bound chain holds the carrier, the
calibrated radius at \emph{its own} equilibrium would have passed; it fails by the widest margin in
the table.

The last row is the one that settles the matter. Everywhere else in this article at least one
parameter was chosen for convenience: $r_c=1.6$ because it was to hand, $a=6.0$ because the first
scans used it. In that row both are determined --- $r_c=1.93$ from the isolated atom's ionisation
energy, and $a=5.5$ from that radius's own equilibrium, located by a five-point scan of the
defect-free chain with a minimum at $5.43$~a$_0$. \emph{At parameters the framework does not get to
choose, no barrier can be quoted.}

\paragraph{Why, in one line of arithmetic.} The totals converge only to first order in the mesh, at
about $0.91$~Ha per unit $h$, so at $h=0.30$ each still carries some $0.27$~Ha of discretisation
error --- a fifth of the total --- while the barrier differenced from the pair is $0.01$~Ha. The
quantity sought is under four per cent of the error in each quantity it is subtracted from, and
would need better than ninety-six per cent cancellation to survive. It does not get it: refining the
mesh alone at fixed everything else moves the value $315\to371\to552\to265$~meV, without converging
and without monotonicity. Brute force cannot close this. Reducing the error merely to the size of
the barrier needs $h\approx0.011$, a $4200^3$ grid; reaching ten per cent of it needs $42000^3$.

\paragraph{Two artifacts, both identified.} What the scans do contain is a monotone drift of
$0.5$--$1.5$~eV per site, the carrier being attracted to a chain end, superposed on a deterministic
wiggle of similar size. The wiggle is grid-locking, and it can be caught in the act: at $h=0.348$
the configurations $c=0.125$ and $c=0.250$ return \emph{identical} energies to five decimals,
because the cell-by-cell assignment did not flip between them. The partition is a staircase pinned
to the grid, and the interface energy jumps as the staircase does. Neither artifact is noise ---
repeated runs of one configuration agree to all five decimals recorded --- which is precisely why
scatter of this kind survived so long undetected: it is exactly reproducible, and reproducibility
was taken for correctness.

\paragraph{A caution about spacing, from the exact side.} Every barrier first computed here used
$a=6.0$. On a hydrogen chain of the same geometry, full configuration interaction puts the binding
at $0.8$~mHa --- $21$~meV for six atoms, against the same atoms separated. \emph{The chain is not
bound at that spacing.} A lattice that does not hold together cannot pin a carrier, so the absence
of a periodic corrugation there is the expected result and not a finding about RealQM. This
invalidates the original comparison with hopping semiconductors on its own, before any question of
numerical resolution arises: those materials are bound solids.

\paragraph{What would be needed instead.} Not a finer grid. The failure is structural to the method
of differencing large totals, so the routes out change what is computed. Either the barrier is
obtained as a path integral of the interface force, $dE/d\lambda=\oint_\Gamma[w]\,v_n\,dS$,
whose error scales with the interface rather than with the volume; or the partition is given
sub-cell resolution so that $E$ becomes a smooth function of where the walls lie instead of a rough
one. Both are projects rather than adjustments, and neither was attempted here.

\paragraph{The control is sound, and a ring shows it.} Lest the failures above be read as a defect
of the specified-partition method itself, the control was checked on a geometry whose symmetry is
exact. A ring of six kernels carrying six electrons, the partition rotated rigidly by $\phi$ in
units of the spacing, must return the same energy at $\phi=1$ as at $\phi=0$ --- the two are related
by a relabelling --- and must be symmetric about $\phi=\tfrac12$. It is, to every digit recorded:
$E(0)=E(1)=0.56101$, $E(0.125)=E(0.875)=0.35258$, $E(0.375)=E(0.625)=-0.21304$, with one pair
disagreeing by $0.011$~Ha and thereby setting the noise floor. So the method reproduces exact
translations exactly, and the chain's asymmetry --- its minimum displaced from $d=\tfrac12$, which a
reflection through a core forbids --- is the finite chain's two end domains, $29\%$ of a seven-core
system, and not the partition. That ring is not otherwise usable: its angular sectors widen with
radius, most of each domain lies far from its kernel, and the total energy comes out
\emph{positive} for a system six hydrogen atoms would bind at $-3$~Ha. It settles the control and
nothing else.

\paragraph{Why this quantity in particular resists measurement.} The construction is a
Frenkel--Kontorova system, and saying so explains the difficulty rather than merely naming it. The
electron domains are a chain with elastic stiffness $\kappa$ supplied by Coulomb repulsion between
neighbours; the kernels are a periodic substrate of amplitude set by $V$. Every configuration in
this article is a displacement field $u(x)$ of electrons against kernels: the rigid slide is its
uniform mode, and a kink is the localised profile carrying $u$ from $0$ to one spacing. In that
model the width of the kink and the barrier to moving it are
\begin{equation}\label{eq:fk}
W=\sqrt{\kappa/K},\qquad K=V''(u_{\min}),\qquad E_{\rm PN}\sim\exp(-c\,W/a),
\end{equation}
so \emph{the barrier is exponentially sensitive to a ratio of two curvatures}. That is the worst
possible form for direct measurement --- an exponential magnifies every error in $\kappa$ and $K$ ---
and it is why differencing total energies was never going to work, whatever the mesh. It also
indicates the route: $\kappa$ and $K$ are local curvatures, taken between configurations differing
infinitesimally, where discretisation errors cancel instead of competing. Measuring those and
inferring $E_{\rm PN}$ from \eqref{eq:fk} is a different calculation from the one attempted here,
and it is left to a separate study. A caution belongs with it: \eqref{eq:fk} assumes
nearest-neighbour coupling, while Coulomb is long-ranged, which changes kink profiles from
exponential to power-law and shifts the transition. It is a guide to the regime, not a formula to
trust to three figures.

\paragraph{What this says about scope.} Less than this article set out to establish, and in a
different direction. No activation energy is available, so the construction is neither placed in
the hopping-semiconductor class nor excluded from it by measurement. The scope claim of
\S\ref{sec:scope} stands or falls on the structural argument instead --- non-overlap, a vanishing
transfer integral, and the absence of any electron equation of motion --- which is where it should
have rested from the beginning, since those follow from the postulate and need no grid at all.

\paragraph{Earlier figures, and why they are not quoted.} A series of barrier scans was made
before the difficulty above was understood, all of them by the free-boundary route: a compression
test on a ring, a doping series at two concentrations, a lattice barrier between interstitial
pockets of order $5$~eV, and a localised interface defect costing $179$~meV independently of chain
length. Every one specified or evolved a partition that no symmetry fixes, so each is a bound of
its stated order rather than a determination, and all are superseded by the barrier above. What
survives from them is the direction of the density trend, which the symmetry-fixed measurement
confirms.

\paragraph{What is actually transported.} The carrier in this framework is not an electron. An
electron here is not a particle but a unit of charge density on a domain, and what exists is the
\emph{partition} of space; dynamics is the evolution of that partition, and transport is the
partition pattern travelling. What moves along the chain is the \textbf{free boundary} --- a locus
where two densities balance --- and nothing material traverses it at all.

The quantitative form of this is the useful one. As a kink passes, each electron advances by exactly
one site spacing and then stops, while the kink advances the entire length of the chain. The
carrier's displacement and its constituents' displacement are decoupled, and their ratio is the
chain length. This is the relation between a dislocation and the slip it produces: the dislocation
crosses the crystal, no atom moves far. Charge is genuinely transported and real charge arrives at
the far end; it is the \emph{carrier} that is immaterial, in the sense that a wave is immaterial
while the water is not.

That diagnosis also accounts for the failures recorded earlier in this study rather than excusing
them. Each of the unsuccessful constructions imposed a carrier that had to \emph{translate}, moving
bodily through occupied space, and returned costs of tens of electron-volts or else collapsed. Those
were the price of forbidding the only mechanism the framework has. The solver states as much
directly: its free boundary evolves by the density balance $(\rho_i - \rho_j)/(\rho_i + \rho_j)$,
flipping ownership cell by cell. Boundary motion is the only motion it performs. The corona result
is the same statement seen from the favourable side --- an added electron spread into a uniform
shell at no cost, which is a boundary redistributing freely.

This raises the stake of the kink measurement. If transport works here it works \emph{because} of
the free boundary, which is the single feature separating this framework from the standard theory;
and if the boundary cannot carry charge cheaply, the difficulty is not in a detail of the
construction but in the framework's defining object.

\paragraph{A repair that is not available, and one that is.} An obvious response to the barrier is
to let valence electrons share a territory, which removes it by construction. The $n^{5/3}$
scaling rules that out: sharing destroys the extensivity the partition supplies, and a metal would
lose its kinetic energy density by fifteen orders of magnitude.

\paragraph{What the framework then is.}\label{sec:hubbard} Non-overlapping sites, occupation $0$, $1$ or $2$, an on-site
repulsion $U$, and transfer between neighbours: that is the \textbf{Hubbard model}. Which locates
RealQM exactly. As formulated --- one unit of charge per domain, always --- it is the Hubbard model in
the limit $U\to\infty$, where double occupancy is forbidden outright and every site is singly
occupied. That limit is a Mott insulator at every density, which is precisely what
the earlier barrier scans found by direct computation, and it explains the density trend as well: no
finite $U$ can be beaten by a bandwidth that does not exist.

The repair is correspondingly minimal. Keep the partition, keep the geometry, keep the Coulomb
energetics, and allow occupation to take the value two at a computed cost. The framework then sits
where the standard treatment of strongly correlated matter sits --- the Hubbard model is the working
description of Mott insulators, of the cuprates, and of the materials for which band theory fails ---
with the difference that its $U$ and its lattice are derived rather than supplied.

We do not carry this out here. What is established is where the obstruction lies: not in the
interface condition, not in the geometry, but in the rule of one unit of charge per domain. That rule
is also the source of one of the framework's results, charge quantization as a consequence of
indivisibility rather than an assumption. The tension is worth stating plainly: \emph{the postulate
that explains why charge comes in units is the postulate that forbids a metal.} Matter has it both
ways, quantizing the particles while leaving the amplitude on each site continuous, and a partition
of real densities has no evident means to do the same.

\section{The carrier itself, against exact diagonalisation}\label{sec:stdqm}

Everything above is internal to the framework. One external check is worth keeping, and the sharpest
form of it does not compare energies at all. Take the same object this article hops --- $n$ protons,
$n-1$ electrons, nuclei clamped --- and diagonalise it exactly in a minimal basis. The Loewdin site
populations then say directly whether the carrier is localised:

\begin{center}
\small
\begin{tabular}{@{}rrrrl@{}}
\toprule
$a$ (a$_0$) & $E_{\rm corr}$ (eV) & $4t$ (eV) & max hole/site & \\
\midrule
$3.0$ & $-6.43$ & $11.11$ & $0.223$ & delocalised\\
$4.0$ & $-11.97$ & $4.65$ & $0.215$ & delocalised\\
$4.5$ & $-14.64$ & $2.96$ & $0.226$ & delocalised\\
$6.0$ & $-19.86$ & $0.71$ & $0.253$ & delocalised\\
\bottomrule
\end{tabular}
\end{center}

\noindent A hole spread evenly over seven sites reads $0.143$; one sitting on a single site reads
$1.000$. At every spacing the ground state puts it on four or five. \emph{There is no localised
carrier, and therefore no barrier for one to cross.} Not because correlation is weak ---
$E_{\rm corr}$ reaches $19.9$~eV --- but because strong correlation and localisation are different
things: a one-dimensional Mott insulator doped with a single hole conducts, the holon propagating
freely.

The disagreement is therefore about the \emph{state} and not about a number, which is the harder
kind to explain away. Two qualifications keep it honest: a minimal basis is exact diagonalisation
within a deliberately crude span --- STO-6G misses the hydrogen atom by six per cent before any
chemistry begins --- and larger bases did not converge for the stretched chains, so no exact
reference exists for these configurations and nothing here declares the pinning wrong for them. What
it does establish is the direction of the error. The construction over-localises, by construction,
being $U\to\infty$: it fails as \S\ref{sec:scope} said it would rather than arbitrarily, and where
the mechanism of this article would actually be tested --- a doped Mott insulator, an organic
semiconductor, a polaronic oxide, where on-site repulsion and not band structure sets the gap --- no
comparison is made here.


\paragraph{Two controls added after the fact, and what each is blind to.} The first is free and
should have been applied from the start: \emph{domains related by a symmetry of the configuration
must have equal eigenvalues}. Each domain's stationarity condition carries an $\varepsilon_i$, and on
a chain where every site has the same environment they must coincide. A total energy is one number
and can look reasonable while the state beneath it is inhomogeneous; the eigenvalues expose that
directly. Applied to the coaxial geometry of the follow-on work it showed terminal domains binding an
electron-volt harder than interior ones --- which identified a corrugation as an end artifact rather
than a registry effect --- and, on longer chains, a smooth bowl across the whole chain that
distinguishes a slowly converging one-dimensional Coulomb sum from a local end effect. It is blind
to exactly what the mesh ladder and the translation control are blind to: an error respecting the
lattice symmetry passes all three, which is how the exclusion rule above survived them.

The second is a way to measure a bulk quantity on a chain that has ends. Displace only the middle $M$
carriers by half a spacing: the energy difference is $M$ times the bulk cost plus two junction costs,
so differencing two window sizes cancels the junctions exactly and no terminal domain moves in either
configuration. The decomposition asserts linearity in $M$, which a third window tests rather than
fits. It is the instrument the rigid slide of a finite chain should have used from the beginning.

\section*{A note on what was tested numerically}

The plasma frequency $\omega_p^2=4\pi n$, the screening relation, the Thomas--Fermi comparison and
the vanishing of $C$ are analytic. The alkali lattice constants are a radial Wigner--Seitz
calculation with the stated pseudopotential and correlation functional, run for five alkalis; the
standard column reproducing the
series is what makes the comparison meaningful, and its own $44\%$ mean error in the bulk modulus
sets the accuracy available at this level of model. Three attempts are reported as null rather than omitted. The interface parameter
$\beta$ was tested on helium at three meshes at fixed imaginary time, giving shifts of $+0.151$,
$-0.018$ and $-0.184$~Ha as the mesh refined --- swinging through zero and growing, with absolute
energies moving away from the exact $-2.9037$. The cause is the geometry rather than the arithmetic:
in helium the two domains meet at a plane \emph{through the nucleus}, so interface error and
Coulomb-cusp error are superimposed and cannot be separated on a uniform grid. The measurement wants
$\mathrm{H_2}$, where the domains meet at the midplane \emph{between} the nuclei in a smooth
region. An interface-condition test on helium was
not converged, giving a sensitivity that reversed sign under mesh refinement, and no conclusion is
drawn from it. Two chain calculations --- five hydrogens with an excess electron, and five hydrogen
kernels with one electron missing --- returned no transport, but for a reason established by reading
the constructions rather than by the runs: a carrier without a kernel feels no lattice force, and a
bare kernel receives no domain for a vacancy to occupy. Both are recorded because the null results
located those obstructions. A caution on the numerics belongs with them. Several attempts returned
energies below the physical floor for their particle content, but these track the boundary speed and
the seeding rather than the configuration: with the interface driven faster than the densities can
follow, or with a carrier seeded directly onto an occupied kernel, the relaxation runs away, while the
same configurations at a quarter of that speed settled to physical values. A two-domain shell on one
nucleus is in fact stable under the density-balance rule --- compressing the inner domain raises its
density at the interface and lowers the outer's, which pushes the boundary back out --- and both
domains carry real kinetic energy there, since the nucleus shapes their profiles. The vanishing of
$C$ is a statement about a \emph{uniform} density and does not transfer to a shaped one. Scripts and pages accompany the source.



\section{Conclusion}\label{sec:conclusion}

The question was whether a construction of non-overlapping unit charge densities, meeting at free
boundaries, conducts. It comes in the two tiers of \S\ref{sec:bar}, and the formulation answers one
of them.

The \textbf{classification} it settles, and settles from a static calculation, which is what Kohn's
result licenses: the polarisation is bounded, the curvature positive, the threshold field beyond
dielectric breakdown, and $\alpha$ extensive in the number of atoms. This is an insulator, and not
by a near miss. What cannot be reached is the \textbf{rate} --- a conductivity, or the activation
energy that would set one --- and the reasons for that are worth more than the number would have
been.

\paragraph{The carrier is a free boundary, and that identification holds.} An electron here is a
unit of charge density on a domain; what exists is the partition of space. Transport is therefore
the partition pattern travelling, in the relation a dislocation bears to slip: charge arrives at the
far end while nothing material crosses the sample. That picture accounts for every null result
reported above --- each had imposed a carrier obliged to translate bodily through occupied space ---
and it is why the one transport calculation this framework performs successfully is Grotthuss
conduction, where the carrier is a proton and the framework does have an equation of motion for it.

The picture has a definite limit, and it should be stated with the claim rather than after it.
Defect-mediated transport is real --- dislocation glide, charge-density-wave sliding and ionic
conduction all proceed that way --- but it is not how a metal conducts. Residual resistivity
\emph{falls} as defects are removed, so a carrier that must migrate as a defect predicts the
opposite of what purity does. The free boundary is therefore the mechanism this construction is
restricted to, not the carrier of ordinary conduction, and its prominence here is a fact about the
formalism rather than a discovery about charge transport.

\paragraph{Costing that motion defeats the method, and the control is what shows it.} A barrier of
$0.01$~Ha differenced from totals of $1.3$~Ha carrying $0.27$~Ha of first-order discretisation error
requires a cancellation it does not get, and brute force cannot supply it: reaching ten per cent of
the barrier needs a $42000^3$ grid. Scanning a full lattice period instead of sampling two
configurations supplies the diagnosis for free, since translation by one spacing must return the
same energy. Six of eight such scans fail that control, the two that pass survive no change of
parameter, and the only scan at parameters the framework does not choose fails by the widest margin.
What the scans do contain is a drift of the carrier toward a chain end and grid-locking of a
staircase interface --- both of order an electron-volt, both exactly reproducible, which is why they
were mistaken for signal for as long as they were.

\emph{But the question does not require that quantity.} Asked as a threshold field rather than as a
barrier --- and on a cyclic chain, where translation by one spacing is an exact relabelling and the
control holds by construction --- it becomes the maximum gradient of a measured periodic curve, and
\S\ref{sec:threshold} gives $7$--$9\times10^{9}$~V/m, twelve orders above the field at which copper
carries an ampere and two to three orders above dielectric breakdown. The same curve yields a
pinning potential of $812$~meV per electron and a positive curvature, hence bounded polarisation and
no DC current. The verdict the article set out to reach is therefore available after all; what was
never available was the number it first tried to reach it with.

\paragraph{Beneath that lies a structural obstacle no refinement touches.} Non-overlap makes the
transfer integral vanish identically, so band conduction is absent in kind: there are no extended
states to carry it. And in this scheme nuclei move under classical equations of motion while
electrons relax to a static minimum --- so there is no electron dynamics, and hence no object in the
theory corresponding to a current. \emph{A minimisation has no time in it.} Every attempt above
infers transport from an energy landscape, which is a proxy, and it is the proxy that breaks. The
missing ingredient is not a better barrier calculation but an equation of motion, and it exists
elsewhere in the programme: the time-dependent form in which a real density carries a complex
current is where a conduction calculation belongs --- and it supplies the current, not its rate,
which needs the dissipation no framework derives unaided.

\paragraph{What the framework is, and is not.} Single occupancy with long-range Coulomb is
\emph{not} the Hubbard atomic limit, whose degeneracy over carrier position would give a barrier of
exactly zero; the Coulomb interaction between extended densities breaks that degeneracy, placing the
construction closer to an extended Hubbard model where charge ordering is genuine physics. Within
bound matter --- structure, binding, geometry --- it works, and the transport regimes it reaches are
those with nuclear carriers. An electronic transport \emph{rate} is outside, and against exact diagonalisation on
the same chains the reason is visible directly: the exact carrier occupies four or five sites where
this construction admits only one.

\paragraph{On the system tested.} Every measurement reported here is on a regular chain of atoms
carrying one valence electron each. The extension to a regular two- or three-dimensional lattice
needs nothing new --- the machinery is three-dimensional already and only the kernel positions
change --- but it has not been made, and a partition with more neighbours to slide against is where
the corrugation would next be tested.

\paragraph{What the time-dependent form would have to supply.} That the framework is static here is
a fact about this article, not about RealQM, and the route out is specific enough to state. What has
to be supplied is a velocity, and it need not arrive as a complex phase: a real density carrying a
real velocity field is Madelung's form of the same equation \cite{madelung}, and in this framework
the natural variables are the densities together with the velocities of the boundaries that bound
them, with $J=\rho v$ and continuity. Either way a current exists as an object, and the first
obstacle goes. It does not remove the second, since a Hamiltonian
dynamics under a constant field rings rather than settles, exactly as the $3N$ equation does, and a
steady current needs dissipation. But the dissipation is already in the scheme and need not be
imported: the nuclei move under classical equations of motion, so energy can leave the electrons
into lattice motion, and the electron--lattice channel that \S\ref{sec:bar} found every treatment
inserting by hand --- a collision term, a partial trace, thermalising reservoirs --- would here be
present in the equations being solved. A resistivity obtained that way would have $\tau$ as an
output. That is the strongest form the claim of this programme could take on transport, and nothing
in this article establishes it. What time-dependence would \emph{not} change is the exclusion:
a phase does not make domains overlap, so there is still no transfer integral and no extended
states, and the target remains activated transport with localised carriers.

\paragraph{The result is therefore a boundary, not a barrier.} A method with a mapped edge is more
usable than one claimed to be universal, and the edge is worth stating exactly rather than
sweepingly, since its four parts are not of the same kind. \emph{Metallic conduction} is excluded
structurally: non-overlap leaves no extended states and no transfer integral, and no refinement
produces them. A \emph{conductivity} is unavailable in the static form, which has no phase, hence no
current as an object and no clock to give one a rate. \emph{Nuclear} transport is inside, protonic
conduction being the one transport calculation that works. And electronic transport by boundary
motion is \emph{described rather than forbidden}: the partition slides, the free boundary carries
the charge, and what this article reports is not the absence of that mechanism but its pinning ---
on one chain, in the rigid mode, on a clamped lattice, with two upper bounds that both push the
number down. A pinned mechanism is a described one. Everything this article originally set out to
measure --- a carrier barrier, a hopping rate --- lay on the far side of the first two.


\begin{thebibliography}{9}
\bibitem{realqm} C.~Johnson, \emph{Real Quantum Mechanics}, \texttt{https://realqm.com}.
\bibitem{realqmbook} C.~Johnson, \emph{Real Quantum Mechanics: a new view of the atom}, book
manuscript.
\bibitem{realnucleus} C.~Johnson, \emph{RealNucleus: nucleons as charge densities}, submitted.
\bibitem{madelung} E.~Madelung, \emph{Quantentheorie in hydrodynamischer Form},
Z.~Phys.~\textbf{40} (1927) 322.
\end{thebibliography}

\end{document}
